
\documentclass[11pt, oneside]{book}
\usepackage{libertine}
\usepackage{datetime}
\usepackage[utf8]{inputenc}
\usepackage[a4paper, total={6in, 9in}]{geometry}
\usepackage{braket}
\usepackage{xcolor}
\usepackage{amsmath}
\usepackage{amsfonts}
\usepackage{amsthm}
\usepackage{amssymb}
%\usepackage[ocgcolorlinks]{hyperref}
\usepackage{hyperref}
%\usepackage{hyperref,xcolor}
%\usepackage[ocgcolorlinks]{ocgx2}
\usepackage{cleveref}
\usepackage{graphicx}
\usepackage{svg}
\usepackage{float}
\usepackage{tikz}
\usetikzlibrary{patterns, shapes.arrows}
\usepackage{adjustbox}
%\usepackage{tikz-network}
\usepackage{tkz-graph}
\usepackage{tkz-berge}
\usepackage[linesnumbered]{algorithm2e}
\usepackage{multicol}
\usepackage[backend=biber,style=alphabetic,sorting=ynt]{biblatex}
%\usepackage{xcolor}
%\usepackage{tkz-berge}
%\usepackage{tkz-graph}
\usepackage{pgfplots}
\usepackage{sagetex}
\usepackage{setspace}
\usepackage{etoc}
%\usepackage{wrapfig}
\usepackage{pgfgantt}
\DeclareUnicodeCharacter{2212}{−}
\usepgfplotslibrary{groupplots,dateplot}
\pgfplotsset{compat=newest}

\newtheorem{theorem}{Theorem}
\newtheorem{definition}{Definition}
\newtheorem{example}{Example}
\newtheorem{claim}{Claim}
\newtheorem{fact}{Fact}
\newtheorem{remark}{Remark}
\newtheorem*{theorem*}{Theorem}
\newtheorem{lemma}{Lemma}
\crefname{lemma}{Lemma}{Lemmas}
\hypersetup{colorlinks=true}
% , allcolors=blue,allbordercolors=blue,pdfborderstyle={0 0 1}}
%\hypersetup{pdfborder={2 2 2}}
% pdfpagemode=FullScreen,
% backref 

\newtheorem{problem}{Problem}
\crefname{problem}{Problem}{Problems}

\DeclareMathOperator{\Ima}{Im}


\usetikzlibrary{positioning}
\addbibresource{./sample.bib} 

%\usepackage[hebrew,english]{babel}
%\titleformat{\chapter}[hang]{\bf\huge}{\thechapter}{2pc}{}

\begin{document}
\newcommand{\commentt}[1]{\textcolor{blue}{ \textbf{[COMMENT]} #1}}
\newcommand{\ctt}[1]{\commentt{#1}}
\newcommand{\prb}[1]{ \mathbf{Pr} \left[ #1 \right]}
\newcommand{\prbm}[2]{ \mathbf{Pr}_{ #2 }\left[ #1 \right]}
\newcommand{\prbc}[3]{ \mathbf{Pr}_{ #2 }\left[ #1 \right | #3]}
\newcommand{\prbcprb}[3]{ \prbc{#2}{#1}{#3} \cdot \prb{#3} } 
\newcommand{\expp}[1]{ \mathbf{E} \left[ {#1} \right]}
\newcommand{\onotation}[1]{\(\mathcal{O} \left( {#1}  \right) \)}
\newcommand{\ona}[1]{\onotation{#1}}
\newcommand{\PSI}{{\ket{\psi}}}
\newcommand{\xij} { X_{ij} } 
\DeclareMathOperator{\Ima}{Im}
%\newcommand{\LESn}{\ket{\psi_n}}
%\newcommand{\LESa}{\ket{\phi_n}}
%\newcommand{\LESs}{\frac{1}{\sqrt{n}}\sum_{i}{\ket{\left(0^{i}10^{n-i}\right)^{n}}}}
%\newcommand{\Hn}{\mathcal{H}_{n}}
%\newcommand{\Ep}{\frac{1}{\sqrt{2^n}}\sum^{2^n}_{x}{ \ket{xx}}}
%\newcommand{\HON}{\ket{\psi_{\text{honest}}}}
%\newcommand{\Lemma}{\paragraph{Lemma.}}
\newcommand{\Cpa}{[n, \rho n, \delta n]}
%\setlength{\columnsep}{0.6cm}
\newcommand{\Jvv}{ \bar{J_{v}} } 
\newcommand{\Cvv}{ \tilde{C_{v}} } 

\newcommand{\Gz}{ G_{z}^{\delta} } 
\newcommand{ \Tann } {  \mathcal{T}\left( G, C_0 \right) }
\newcommand{\ireducable}{ireducable \hyperref[ire]{[\ref{ire}]} }
\newcommand{\cutUU}{E(U_{-1} \bigcup U_{+1} ,U)} 
\newcommand{\wcutUU}{w\left( E(U_{-1} \bigcup U_{+1} ,U)  \right)}
\newcommand{\testgo}{  \mathcal{T}\left(J, q , C_{0}\right) } 

\newcommand{\duC}{\left( C_{A}^{\perp}\otimes C_{B}^{\perp} \right)^{\perp}}
\newcommand{\duduC}{\left( C_{A}\otimes C_{B}\right)^{\perp}}
  




\begin{titlepage}
    \begin{center}
        \vspace*{1cm}
        
        \includegraphics[width=0.15\textwidth]{huji_logo_notext.pdf}\\
        { \large The Hebrew University of Jerusalem\\
        The Rachel and Selim Benin School of Computer Science and Engineering }
        
        \vspace{2cm}
        
        {\huge \textbf{Ph.D Research Aims.}}
        
        \vspace{1cm}
       % \selectlanguage{hebrew}
       % {\Large} % שם העבודה}
       % \selectlanguage{english}
        
        \vspace{1.5cm}
        
        { \large David Ponarovsky }
        
        \vspace{1cm}
        
        { \large  Under the supervision of Prof. Michael Ben Or }

        
        \vfill
        
        {\large August 2023 }
    \end{center}
\end{titlepage}

\setlength{\parindent}{0pt}
\setlength{\parskip}{5pt}

\chapter{Research Aims.}

Although it may seem that quantum computing is just around the corner, we still do not understand it as well as we understand classical computing. The purpose of my Ph.D. studies is to bridge this gap by focusing on the quantum PCP conjecture, which is known to be true in the classical case and is believed to hold in the quantum regime as well. The conjecture is stated as follows: consider the promise decision problem, where a classical description of a local Hamiltonian is provided. One must determine whether the ground energy is greater than $b$ or less than $a$, while being promised that $b-a$ is greater than a constant number independent of the description length. This problem is \textbf{QMA}-complete, which means that any decision problem whose solution can be verified by a quantum machine in a reasonable time can also be reduced to a decision problem in which, using a quantum machine, one can verify a candidate solution by doing a constant amount of work and failing in probability proportional to the closeness to an actual valid proof.

As the proof for the classical conjecture relies on the existence of error correction codes with specific properties, our strategy is to enhance our understanding of quantum error correction codes. In particular, we aim to develop quantum testable codes, which are codes that allow us to determine whether a given state is in the code by querying a constant amount of bits and failing with a probability proportional to the distance of the state from the code. 





\end{document}
